\documentclass[11pt]{article}

\usepackage[margin=1in]{geometry}
\usepackage{amsmath,amssymb,amsthm}
\usepackage{booktabs}
\usepackage{xcolor}
\usepackage[colorlinks=true,linkcolor=blue!60!black]{hyperref}
\usepackage{microtype}
\usepackage{enumitem}

\emergencystretch=3em

\newcommand{\E}{\mathbb{E}}
\newcommand{\sg}{\mathrm{sg}}
\newcommand{\R}{\mathbb{R}}
\newcommand{\numgood}[1]{\textbf{\textcolor{green!45!black}{#1}}}
\newcommand{\numbad}[1]{\textbf{\textcolor{red!65!black}{#1}}}

\newtheorem{proposition}{Proposition}
\newtheorem{lemma}{Lemma}
\newtheorem{corollary}{Corollary}
\theoremstyle{definition}
\newtheorem{claim}{Claim}
\newtheorem{prediction}{Prediction}
\theoremstyle{remark}
\newtheorem{remark}{Remark}

\title{The Phenomenon:\\ Why the Bilinear Form Works\\
\large ogbench-sota --- phenomenon note, 2026-07-20\\
\normalsize (companion: \texttt{PAPER/rescuing-fb/phenomenon.tex},
\texttt{PAPER/rescuing-fb/understanding-bilinear.tex})}
\author{Factored-FB project}
\date{}

\begin{document}
\maketitle

\section{The question, and the facts any answer must reproduce}\label{sec:facts}

Swapping a monolithic MLP head $Q_\theta([s,a,g])$ for the normalized inner product
$Q(s,a,g)=\phi(s,a)^\top\psi(g)/\sqrt d$ inside GC-IQL turns 0.71 into
\numgood{0.90} on cube-single, transfers to GCIVL ($0.50\to0.79$) and HIQL
($0.15\to0.41$, once routed around the shared subgoal rep), and the same form recurs in
essentially every method that has ever worked on task families: SF/USFA, FB, CRL, VIP, BVN,
HILP/METRA, quasimetric nets. \emph{Why?} Any explanation must reproduce all ten facts:

\begin{center}\small
\renewcommand{\arraystretch}{1.22}
\begin{tabular}{@{}l p{10.4cm} l@{}}
\toprule
 & Fact & Source \\
\midrule
F1 & Bilinear $0.904$ vs MLP $0.713$; the split is the \emph{goal} axis
($\Delta_{\mathrm{GOAL}}$: bilinear $\in[-0.08,+0.08]$, MLP $-16$ to $-24$) & OOD grid \\
F2 & The MLP is the \emph{better} critic fully in-distribution (C1: $0.98$ vs $0.87$) & OOD
grid \\
F3 & The bilinear critic is \emph{not} a better value function: counterfactual rank-AUC MLP
$0.81$ vs bilinear $0.39$ (below chance); the MLP's signal is all in the action-independent
baseline, the bilinear's baseline is anti-calibrated & cf-probe \\
F4 & The bilinear critic's product is the \emph{on-support action ordering}: frozen-critic
reranking $0.91$--$0.96$ ($=$ its trained actor) vs MLP $0.57$--$0.59$ ($<$ its actor); FB's
ordering is anti-informative (declines with $N$: $0.52\to0.17$) & rerank \\
F5 & Action-dependence is $<1\%$ of $Q$-variance for \emph{every} critic (bilinear
$0.46\%$, FB $0.21\%$, MLP $0.10\%$) --- the loss barely sees the thing that decides success
& variance probe \\
F6 & Rank is a floor, not a dial: $d{=}5$ collapses ($0.13$), $d\in[28,512]$ flat
($0.86$--$0.93$); at $d{=}512$ the ``bottleneck'' is as wide as the trunk and still beats the
MLP by 16 & rank sweep \\
F7 & Eval-time goal smoothing cannot create the MLP's missing signal (monotone
\emph{decline}); if pooling is the mechanism it acts at training time & T1 \\
F8 & $\psi(g)\to\psi(s,g)$ destabilizes outright: $2/3$ seeds diverge via a norm runaway on
the state-conditioned goal leg; the stable seed lands at MLP level & bilinear\_sg, T2 \\
F9 & Every \emph{explicit} vector/feature-TD estimator of the same structure collapses
($0.00$--$0.38$); scalar TD through the same structure reaches $0.90$ & SF ladder \\
F10 & The wins and failures are consumption-dependent and domain-dependent: cube
flow $0.90$ / DDPG+BC $0.51$ (off-support peak, frac $0.639$); antmaze DDPG+BC $0.837$ (no
peak, frac $0.003$); scene loses with both actors (no peak --- different failure) & T3,
extraction \\
\bottomrule
\end{tabular}
\end{center}

No single-clause story survives this table: ``better values'' dies on F3, ``compression /
regularization'' dies on F6, ``eval-time denoising'' dies on F7, ``data coverage'' died on
T3's dimension-corrected control. The explanation is a conjunction of four layers, each
independently evidenced, and it is the conjunction --- not any layer alone --- that produces
the cube result.

\section{Layer 0 (representation): policy evaluation is genuinely bilinear}
\label{sec:layer0}

\begin{proposition}[Task-linearity of evaluation]\label{prop:linearity}
For any fixed policy $\pi$, the map $r\mapsto Q^\pi_r$ is linear:
$Q^\pi_r(s,a)=\int M^\pi(s,a,\mathrm ds')\,r(s')$ with $M^\pi$ the successor measure. Hence
if a task family's rewards lie ($\varepsilon$-approximately) in a $d$-dimensional span,
$r_g(s')\approx\beta(s')^\top w(g)$, then \emph{exactly}
\begin{equation}
Q^\pi_g(s,a)\;\approx\;\underbrace{\Big(\textstyle\int M^\pi(s,a,\mathrm
ds')\beta(s')\Big)}_{=:\,\phi^\pi(s,a)\ \text{(the SF of }\beta\text{)}}{}^{\top}\,w(g),
\label{eq:sf}
\end{equation}
i.e.\ the bilinear head is not an approximation \emph{ansatz} but the \emph{canonical form}
of policy evaluation restricted to a low-rank reward family. The goal family of cube-single
--- indicator balls $r_g=\mathbf 1[\|c(s)-g\|\le\varepsilon]$ with $g$ on a smooth
3-manifold --- has a Gram operator with fast spectral decay (a smooth kernel on a 3-dim
index set), so small $d$ suffices and the required $d$ is set by the task family's intrinsic
complexity, not by state dimension.
\end{proposition}

This single proposition already fixes three facts. \textbf{F6}: $d$ must exceed the
$\varepsilon$-rank of the reward family (the $d{=}5$ cliff = cannot span; the flat band =
excess rank is unused because the target itself is low-rank --- consistent with the empirical
near-low-rankness of control $Q$-matrices, Yang et al.\ 2020; Shah et al.\ 2020).
\textbf{F1}: for every goal, $Q(\cdot,\cdot,g)\in\mathrm{span}\{\phi_1,\dots,\phi_d\}$ ---
a \emph{fixed} $d$-dimensional function space fitted against \emph{all} training goals --- so
a novel evaluation goal requires extrapolating only the small map $\psi:\R^3\to\R^d$, never
the value function over $(s,a)$; the MLP, nonlinear in $g$, must extrapolate the value
itself and pays $16$--$24$ points. \textbf{F9} (half): the structure being right says
nothing about how to estimate it --- see Layer 2.

\begin{remark}[Value equivalence: the bilinear head spends its capacity exactly on the value
span]\label{rem:ve}
The value-equivalence principle (Grimm et al.\ 2020, 2021) says a model/representation need
only be correct on the span of value functions one will ever evaluate. Equation
\eqref{eq:sf} is the value-equivalent compression of the environment for the task family:
$\phi$ retains precisely the $d$ occupancy moments the family's values need, and nothing
else. The expectation-model line (Wan et al.\ 2019, 2021) supplies the closure property that
makes TD work here: when values are linear in features, the Bellman backup needs only the
\emph{expected next feature} --- a linear-in-$\phi$ value family is closed under expectation
backups, so scalar TD on the head is a well-posed operator on the span. A monolithic MLP has
no such closed span; each goal's backup is free to bend the function elsewhere.
\end{remark}

\begin{remark}[The caveat that produces the residual]\label{rem:quasimetric}
The \emph{optimal} value is not linear in $r$ (the $\max$ breaks linearity), and
$-V^*$ of goal-reaching is a quasimetric that symmetric inner products cannot represent with
low distortion (QRL, IQE, MRN). So Proposition~\ref{prop:linearity} can only hold
approximately for the trained (expectile-improved) object, and the approximation error must
land somewhere. Our account: the in-sample expectile keeps the trained object near the
behavior-supported evaluation envelope where linearity is good, and the irreducible
asymmetric residual is pushed \emph{off-support} --- observed as the off-support peak
(\textbf{F10}: cube frac $0.639$), whose location is a property of the domain's value
geometry (antmaze $0.003$), and as the scene failure, where the task family is
combinatorial (multi-object) and a single rank-$d$ product of one state tower and one goal
tower is the wrong span. This is the one honest boundary of the explanation, and it is
\emph{predictive}, not hand-wavy: see Prediction~\ref{pred:peak}.
\end{remark}

\paragraph{Why the same form recurs across a decade of methods.} Every historically
successful task-family method is a rank factorization of the same evaluation functional,
differing only in estimator:

\begin{center}\footnotesize
\renewcommand{\arraystretch}{1.25}
\setlength{\tabcolsep}{4pt}
\begin{tabular}{@{}l l l l@{}}
\toprule
Method & factorized object & estimator & weighting \\
\midrule
SF/USFA & $Q^\pi_w=\psi^\pi(s,a)^\top w$ & feature-TD $+$ regression & task \\
FB & $M^{\pi_z}\approx F^\top B\rho$ & density (measure) TD & occupancy $\rho\otimes\rho$ \\
CRL & $\log\frac{p^{\pi}(s^+|s,a)}{p(s^+)}\approx\phi^\top\psi$ & InfoNCE on positives
$\sim M$ & trajectory-future \\
VIP & first-passage dual value $-\|\varphi(s){-}\varphi(g)\|$ & Fenchel-dual
time-contrastive & goal-anchored \\
BVN & $Q(s,a,g)=f(s,a)^\top\phi(s,g)$ & scalar TD (online, HER) & hindsight \\
GC-IQL bilinear (ours) & $Q=\phi(s,a)^\top\psi(g)/\sqrt d$ & scalar expectile TD &
hindsight mix \\
\bottomrule
\end{tabular}
\end{center}

\noindent The form keeps working because the \emph{object is genuinely bilinear}
(Proposition~\ref{prop:linearity}); when one of these methods fails --- FB on manipulation
--- the bridge program showed by exhaustive one-knob morphing that the failure is the
\emph{estimator column}, never the form (identical architecture: \numgood{87} under scalar
TD, \numbad{16} under the measure; the companion note carries that phenomenon). ``Bilinear
works'' and ``FB fails'' are the same finding viewed from two sides: \textbf{structure
$\neq$ estimator.} Our configuration is, in taxonomy terms, \emph{successor features whose
reward basis $\psi(g)$ is co-learned by scalar hindsight TD} --- SFs without feature-TD.

\section{Layer 1 (statistics): the shared basis pools a $<\!1\%$ signal across goals}
\label{sec:layer1}

Layer 0 says the form is right; it does not yet explain F3/F4 (ordering vs calibration) ---
why the bilinear critic \emph{finds} an action ordering the MLP misses on the \emph{same}
data and loss.

\begin{proposition}[Gradient coupling / cross-goal pooling]\label{prop:pooling}
Decompose $Q(s,a,g)=U(s,g)+A(s,a,g)$ and the advantage into goal-common and
goal-idiosyncratic parts, $A(s,a,g)=\bar A(s,a)+\epsilon(s,a,g)$ with
$\E_g[\epsilon]=0$. Under the bilinear head the TD gradient at $\phi(s,a)$ is
\begin{equation}
\nabla_{\phi(s,a)}\mathcal L=\tfrac{1}{\sqrt d}\,\E_{g}\big[\delta(s,a,g)\,\psi(g)\big],
\label{eq:pool}
\end{equation}
a \emph{linear} pooling of the TD errors of every relabelled goal along their
$\psi$-directions: goal-common advantage errors accumulate coherently over the goal
distribution while zero-mean idiosyncratic errors cancel at rate
$O(N_g^{-1/2})$. The shared tower is a multi-task regression with a shared representation;
its effective sample size for the goal-common component $\bar A$ is (state-visits
$\times$ goals), not per-goal. A monolithic MLP receives only per-input gradients; with
$\mathrm{Var}[A]<10^{-2}\,\mathrm{Var}[Q]$ (F5), the variance-dominant baseline owns its
capacity and optimization signal (and TD's implicit rank collapse pushes the same way), so
its action signal is \emph{absent}, not noisy --- which is exactly what F7 measured. On cube
the goal-common component is large by construction: post-grasp, ``which action keeps the
grasp while progressing'' is nearly the same ordering problem for every transport goal (the
phase diagnostic: the entire method spread is post-grasp,
$P(\mathrm{success}\mid\mathrm{secured})$ $0.98$ vs $0.80$).
\end{proposition}

This layer fixes \textbf{F3--F5, F7}, including the capacity-allocation \emph{trade}: the
rank-$d$ head cannot simultaneously match an arbitrary baseline surface $U(s,g)$ (itself a
rank-$d$-constrained matrix over states $\times$ goals) and the ordering; it spends its
freedom on the pooled, goal-consistent structure and pays with calibration
(anti-calibrated $U$, F3) --- which is downstream-irrelevant, because extraction consumes
orderings (Layer 3). The MLP makes the opposite trade: calibrated baseline (cf-AUC $0.81$),
no ordering.

\begin{prediction}[T5, the discriminating control --- pre-registered]\label{pred:t5}
The pooling mechanism lives in the \emph{linear} readout: only an inner product makes
$\nabla_\phi\mathcal L$ a $\psi(g)$-weighted linear pool \eqref{eq:pool}. A
\emph{late-fusion} control (identical two towers, MLP readout on $[\phi;\psi]$) re-mixes
nonlinearly per input and breaks the coherent pooling. Prediction: late-fusion reranks near
MLP level ($\le0.65$), far below bilinear ($0.91$--$0.96$), despite identical towers and
parameters. A parameter-matched MLP (the $2\times$ trunk confound control) stays at MLP
level. If late-fusion instead reranks at bilinear level, Layer 1 is falsified and the
explanation reduces to Layers 0/2 --- this is the single most informative cheap experiment
on the board.
\end{prediction}

\section{Layer 2 (optimization): why scalar TD trains this structure and nothing else
does}\label{sec:layer2}

\textbf{(a) Scale and grounding.} Scalar TD pins units ($|Q|\le(1-\gamma)^{-1}$, reward
units); the $\sqrt d$ normalization keeps the product's scale $O(\|\phi\|\|\psi\|)$-stable
--- the same principle as bounded/structured critic outputs (HL-Gauss, SimbaV2). Vector
feature-TD and density-TD are self-referential: their fixed points are invariant to
invertible linear reparameterizations, admit collapse ($F{=}0$-type) solutions unless
anchored, and their anchors import density units with unbounded dynamic range --- the
measured runaways ($q\to10^{13}$; SG-ISFB inflation) and the F9 ladder ($0.00$--$0.38$
across every explicit vector-TD estimator of the very same structure) are one defect.
\emph{Scalar TD through the structure} is the unique tested estimator that both grounds the
scale and lets the structure emerge.

\textbf{(b) The $\psi(g)$ bottleneck as global coupling (F8).} With $\psi$ a function of $g$
alone, any local score inflation at $(s,g)$ must move \emph{globally shared} features and is
corrected by the rest of the data; $\psi(s,g)$ opens a local $(s,g)$ pocket that expectile
TD's upward pressure (no output bound, no $V$-target on that leg) inflates without
opposition --- T2 located the divergence precisely there ($\E\|\psi\|^2$
$1083\to2355\to3789$ monotone in severity, $\phi$ sane throughout). The same coupling is
what lets goal-shift generalization be free (F1): the head cannot even \emph{represent} a
goal-idiosyncratic exception. Note the regime inversion vs online BVN (which prefers
$f(s,a)^\top\phi(s,g)$): the offline expectile backup with no output bound is the standing
suspect, testable by bounding the $\psi(s,g)$ leg (implementation note \S\ref{sec:impl}).

\section{Layer 3 (consumption): the product is an ordering, and the consumer must be
anchored}\label{sec:layer3}

The trained bilinear critic's asset is the on-support ordering (F4); its one structured flaw
is the off-support maximum (Remark~\ref{rem:quasimetric}). Extraction decides which of the
two you consume: BC-anchored flow / selection-by-argmax consumes exactly the ordering and
never visits the flaw ($0.90$--$0.96$); gradient ascent (DDPG+BC) walks distance
$17.9$ (whitened) off-manifold to the fabricated peak ($0.51$), while the MLP --- no peak,
no ordering --- is indifferent ($\approx0.73$ either way). Domains flip the sign because the
peak's existence is domain value-geometry (F10). The corollary design rules, both verified:
(i) on manipulation data the BC anchor is \emph{mandatory} (the support argument in the
companion note: the stable-vs-slipping boundary has $\rho$-mass $1.7\times10^{-5}$ --- no
critic may be asked to rank stability counterfactuals); (ii) \emph{the bilinear head must
not train a representation the policy shares} (HIQL: routing the bilinear TD loss around
$\zeta$ flips $14.5\to41.3$).

\section{The explanation, assembled}\label{sec:assembled}

\begin{claim}[Why the bilinear form works]\label{claim:final}
$Q(s,a,g)=\phi(s,a)^\top\psi(g)/\sqrt d$ trained by scalar hindsight expectile TD works
because four independently-evidenced layers align:
\textbf{(0)} the target object --- policy evaluation over a low-complexity task family ---
is \emph{itself} bilinear (Prop.~\ref{prop:linearity}), so the head is the value-equivalent
parameterization, rank cost is a floor set by the family (F6), and goal generalization is
linear-span reuse, not extrapolation (F1);
\textbf{(1)} the linear readout turns TD into a cross-goal gradient pool
(Prop.~\ref{prop:pooling}) that accumulates the goal-common, $<1\%$-variance action
ordering the task actually needs, at the price of per-goal calibration it does not need
(F3--F5, F7);
\textbf{(2)} scalar TD is the unique tested estimator that grounds this structure's scale
and reward-relevant subspace --- vector/density TD on the identical structure collapses
(F9) --- and the $g$-only $\psi$ leg globally couples the fit, closing the local-inflation
pockets that destroy $\psi(s,g)$ (F8);
\textbf{(3)} anchored/selection extraction consumes exactly the ordering the head produces
and never visits the quasimetric residual it necessarily exiles off-support
(F10). The historical record (SF, FB, CRL, VIP, BVN) is layer 0 recurring under different
estimators; the FB-on-manipulation failure is a layer-2/weighting failure of one estimator,
not of the form. Where any layer's precondition fails, the form fails exactly as predicted:
rank below the family's span ($d{=}5$), a compositional family outside one product's span
(scene), gradient consumption of the residual (cube DDPG+BC), or a shared representation
trained by the head (HIQL form a).
\end{claim}

What makes this account hard to challenge is not rhetoric but coverage and falsifiability:
every clause carries a completed experiment, the rival single-clause stories are each dead
on a named fact, and the two genuinely open sub-questions have pre-registered discriminating
experiments (Prediction~\ref{pred:t5}; Prediction~\ref{pred:peak}) whose either outcome
\emph{refines} the layer structure rather than overturning F1--F10, which are measurements.

\begin{prediction}[Peak geometry --- pre-registered]\label{pred:peak}
If the off-support peak is the exiled quasimetric residual, its location should track
embedding-norm growth off-manifold: $\|\phi(s,a)\|$ along rays
$a_{\mathrm{ref}}\to a_{\mathrm{DDPG}}$ grows on cube and not on antmaze, and a norm-capped
$\phi$ (spectral or hard cap) shrinks the cube peak fraction from $0.639$ toward the antmaze
$0.003$ while leaving flow-actor success within noise. If the peak survives norm capping,
the residual-relocation reading is wrong and the peak needs a new account --- but the
extraction rule (Layer 3) stands on its own measurements either way.
\end{prediction}

\section{Final section: proposed experiments and implementations}\label{sec:impl}

Ordered by information-per-GPU-hour; each is specified concretely in
\texttt{implementations.md} (this folder).

\begin{enumerate}[leftmargin=1.7em,itemsep=3pt]
\item \textbf{T5 late-fusion + param-matched MLP} (Prediction~\ref{pred:t5}). Two towers
identical to bilinear; readout MLP$([\phi;\psi])$; rerank protocol at $N{=}32$ plus full
training run; second arm: MLP with $2.06\times$ width to close the parameter confound.
Decides Layer 1.
\item \textbf{Norm-capped $\phi$ / peak-geometry probe} (Prediction~\ref{pred:peak}).
Recompute the T3 peak statistic along reference$\to$DDPG rays vs $\|\phi\|$; then retrain
cube DDPG+BC with a capped $\phi$. Decides the residual-relocation reading and, if
positive, yields a fix that makes gradient extraction safe.
\item \textbf{srank/effective-rank trajectories} of MLP vs bilinear $\phi$ during training
--- the implicit-under-parameterization half of Layer 1's ``absent, not noisy''; cheap, uses
existing checkpoints.
\item \textbf{T9: rank of the true advantage.} Estimate
$\mathrm{rank}_\varepsilon(A^*(s,\cdot,g))$ over the goal grid from counterfactual rollouts
(atlas banks + simulator). Tests Layer 0's premise that the task family's ordering demands
are genuinely low-rank; calibrates the $d$ floor against ground truth.
\item \textbf{Bounded $\psi(s,g)$} ($\tanh$ output bound or $V$-target on the goal leg).
Tests whether F8's divergence is the only obstruction to the BVN pairing offline --- if it
stabilizes at bilinear level, the online/offline inversion is an optimization artifact, and
the paper's claim narrows cleanly to ``$g$-only $\psi$ is \emph{sufficient} and \emph{safe},
state-conditioning is fixable but unnecessary.''
\item \textbf{Scene split-head}: per-factor goal towers $\sum_k\phi_k(s)^\top\psi_k(g_k)$
(Schaul-style factored family). Tests Remark~\ref{rem:quasimetric}'s account of scene as a
span failure of one product; a positive result extends Layer 0 to compositional families and
likely recovers the scene loss.
\item \textbf{Cross-pollination (the bridge back to research question 2).} Layer 1 says the
mechanism is \emph{training-time} pooling through a shared basis under scalar TD; the
companion note's RG-FB v1 (task-prior scalar grounding) is exactly ``make the pooling
distribution richer than Dirac goals.'' The two programs share one sentence: \emph{the
bilinear form is the right structure because evaluation is bilinear; scalar task-weighted TD
is the right estimator because it is the only one that grounds, pools, and stays on
support.}
\end{enumerate}

\end{document}
